\begin{pastexam}[label={px:ntu-114-dsa-14}]{臺大資網所 114 學年度 第 14 題}
（複選；每個選項單獨計分，答對 1 分，答錯 0 分）

(5 points) Let $G$ be an $n\times n$ grid graph, with each undirected edge having a certain weight.
Specifically, the vertices of $G$ can be labeled as $(i,j)$ where $1\le i,j\le n$. There is an edge between
$(i,j)$ and $(k,\ell)$ if and only if $|i-k|+|j-\ell|=1$. Suppose we want to compute the \textit{value} of a
minimum cut that separates $(1,1)$ and $(n,n)$. The value of a cut is defined as the sum of weights to all edges
crossing the cut.

Four students proposed different ideas for solving this task. Let us examine their approaches and select all
correct statement(s) below.
\begin{enumerate}[label=(\Alph*)]
\item Student A: Some minimum cut must cut through at most $2n$ edges. Since there are $2n(n-1)\le 2n^2$ edges in
$G$, we can design a brute force algorithm that enumerates at most $\sum_{k=1}^{2n}\binom{2n^2}{k}$ subsets of
edges and check: if the removal of those edges separates $(1,1)$ and $(n,n)$. Any such subset leads to a
candidate for the minimum cut.
\item Student B: For some minimum cut, the crossing edges must form a monotonic path from $(n,1)$ to $(1,n)$. In
other words, those edges form a path of exactly $2n-2$ edges from $(n,1)$ to $(1,n)$. Thus, we can solve this
problem via dynamic programming.
\item Student C: Why don't we simply computes a maximum flow and then turning it to a minimum cut? First,
construct a flow network by letting $(1,1)$ be the source and $(n,n)$ be the sink. For each undirected edge
connecting $(i,j)$ and $(k,\ell)$, create two directed edges in opposite directions, both with capacities equal
to that edge's weight. By using Dinitz's (Dinic's) algorithm, the maximum flow from $(1,1)$ to $(n,n)$ can be
computed in $O(n^6)$. Using the Max-Flow Min-Cut Theorem, the returned maximum flow value equals the value of
the minimum cut.
\item Student D: To separate $(1,1)$ and $(n,n)$, for some minimum cut, the crossing edges must form a shortest
path from the south-west border $WS:=\{(i,1)\mid 1<i\le n\}\cup\{(n,j)\mid 1\le j<n\}$ to the north-east border
$NE:=\{(i,n)\mid 1\le i<n\}\cup\{(1,j)\mid 1<j\le n\}$. Thus, we can run the Bellman-Ford algorithm that seeks for
any shortest path from a vertex in $WS$ to any vertex in $NE$. Finally, the length of the shortest path is the
value of the minimum cut.
\item None of the above choices.
\end{enumerate}
\end{pastexam}

\begin{pastsoln}
答案：(C)。

\textbf{背景}。 網格是平面圖。 把每個格子（面）當成一個點、外面的區域沿 $(1,1)$ 和 $(n,n)$ 切成兩塊各當一個點，相鄰的面之間
連一條「對偶邊」，權重等於它跨過的那條原邊。 分開 $(1,1)$ 和 $(n,n)$ 的割，恰好對應對偶圖中從一塊外部區域走到另一塊的
路徑；最小割的值就是對偶圖的最短路徑長。 這是正確的快速解法，但它是在\textbf{對偶圖}上找最短路徑，路徑可以任意彎曲。

\textbf{讓最小割很長的例子}。 選一條在格子之間蛇行的對偶路徑：從左邊界進入左下角的格子，沿最下一排格子往右，上一排往左，
再上一排往右……走遍全部 $(n-1)^2$ 個格子，最後從上邊界出去。 它跨過的 $(n-1)^2 + 1$ 條原邊權重設 $1$，其他邊權重設
$n^4$。 用到任何一條重邊的割都超過 $n^4$；只用輕邊的割必須包含整條蛇行路徑跨過的邊（缺一條就會留下一條連通
$(1,1)$ 和 $(n,n)$ 的通道），所以最小割唯一，值是 $(n-1)^2 + 1$，用了 $(n-1)^2 + 1$ 條邊。 例如 $n = 5$ 時最小割是
$17$ 條邊，而 $2n = 10$（用最大流驗算，$n = 4, \ldots, 7$ 都是這個結果）。

\begin{itemize}
\item[(A)] 錯。 前提「一定有某個最小割只用到至多 $2n$ 條邊」不成立：上面的例子 $n \ge 4$ 時 $(n-1)^2 + 1 > 2n$，最小割
只有一個，邊數超過 $2n$。 所以只列舉至多 $2n$ 條邊的子集合會漏掉最小割。 （而且就算前提成立，$\sum_{k \le 2n}
\binom{2n^2}{k}$ 也是指數級的。）
\item[(B)] 錯。 同一個例子：最小割的邊不是長度 $2n - 2$ 的單調路徑。 割邊在對偶圖上是一條可以來回彎曲的路徑。
\item[(C)] 對。 無向邊換成兩條反向、容量相同的有向邊，最大流等於無向圖的最小割（資結書定理~26.6，無向圖的處理見
資結書第~26.1.2~節）。 Dinic 演算法的時間是 $O(V^2 E)$；這裡 $V = n^2$、$E = 2 \cdot 2n(n-1) = O(n^2)$，所以是
$O(n^4 \cdot n^2) = O(n^6)$。 敘述完全正確。
\item[(D)] 錯。 方向對了一半：最小割確實對應「從一邊界到另一邊界的最短路徑」，但那是\textbf{對偶圖}上的路徑；在原圖 $G$ 的
頂點上跑 Bellman--Ford 算的是另一回事。 反例 $n = 2$：$WS = \{(2,1)\}$、$NE = \{(1,2)\}$。 四條邊權重：
$(1,1)$--$(1,2)$ 是 $1$、$(2,1)$--$(2,2)$ 是 $1$、另外兩條是 $10$。 割掉兩條權重 $1$ 的邊就分開了 $(1,1)$ 和 $(2,2)$，
最小割是 $2$；但 $G$ 上 $(2,1)$ 到 $(1,2)$ 的最短路徑，兩條路都要經過一條權重 $10$ 的邊，長度 $11$。
\item[(E)] 錯，(C) 是對的。
\end{itemize}
\end{pastsoln}
